% =====================================================================
%  capitulos/02_datos.tex  --  SOLO CONTENIDO
% =====================================================================
%  Sección 2 del informe. Aquí practicamos: etiquetas y referencias
%  cruzadas, figuras, tablas y cajas de aviso.
% =====================================================================
\section{Datos enlazados a gran escala}
% \label da un "nombre" a esta sección para poder citarla desde otro
% sitio con \ref{sec:datos}. Convención: prefijo + tema
% (sec: sección, fig: figura, tab: tabla).
\label{sec:datos}

El formato \emph{HDT} divide un conjunto de datos RDF en tres
componentes: cabecera, diccionario y tripletas~\cite{fernandez2013hdt}.
Esa separación es, en el fondo, la misma idea que aplicamos hoy a
LaTeX: \textbf{dividir en módulos con una responsabilidad clara}.

% ---- FIGURA ----
% Un entorno "figure" crea una figura numerada con pie y etiqueta.
%   [H]         = colócala exactamente aquí (lo permite el paquete float)
%   \centering  = centrar
%   \includegraphics[width=...]{fichero}: inserta la imagen.
%        width=0.85\linewidth = 85 % del ancho del texto
%        El fichero se busca en Recursos/ gracias a \graphicspath (.sty).
%        Se usa el PDF vectorial (nítido a cualquier zoom), no el PNG.
%   \caption    = texto del pie ("Figura 1: ...")
%   \label      = nombre para citarla con \ref{fig:hdt}
% PARA CAMBIAR LA FIGURA: guarda otra en Recursos/ y cambia el nombre.
\begin{figure}[H]
  \centering
  \includegraphics[width=0.85\linewidth]{hdt_esquema.pdf}
  \caption{Esquema de HDT (figura vectorial).}
  \label{fig:hdt}
\end{figure}

% ---- TABLA ----
% tabular define las columnas con {lll}: tres columnas alineadas a la
% izquierda (c = centrada, r = derecha).
%   &   separa columnas     \\   termina una fila
%   \toprule \midrule \bottomrule: líneas horizontales de booktabs
%   @{}...@{} en la definición quita el margen lateral sobrante.
% PARA AÑADIR UNA FILA: copia una fila existente y cambia su texto
% (no olvides los & ni el \\ final).
\begin{table}[H]
  \centering
  \begin{tabular}{@{}lll@{}}
    \toprule
    \textbf{Componente} & \textbf{Contenido} & \textbf{Función} \\
    \midrule
    Header     & Metadatos        & Describe el conjunto de datos \\
    Dictionary & Cadenas de texto & Evita repetirlas \\
    Triples    & Estructura       & Representa el grafo \\
    \bottomrule
  \end{tabular}
  \caption{Los tres componentes de HDT.}
  \label{tab:hdt}
\end{table}

% \ref{fig:hdt} y \ref{tab:hdt} se sustituyen por los números reales.
La figura~\ref{fig:hdt} y la tabla~\ref{tab:hdt} resumen la idea.

% El mismo entorno "aviso", ahora con título propio (argumento opcional
% entre corchetes).
\begin{aviso}[Regla práctica]
  Cuando copies y pegues un bloque por tercera vez, es hora de
  convertirlo en un entorno.
\end{aviso}
